\chapter{Multi-CLD Validation of Context-Insensitive Hallucination Metrics (RQ2)}
\label{ch:rq2_results}

\section{Dataset Coverage}

This analysis evaluates context-insensitive (CI) metrics for hallucination detection across multiple Causal Loop Diagrams (CLDs) within the health domain. The validation dataset comprises:

\begin{itemize}
    \item \textbf{Three CLDs:}
    \begin{enumerate}
        \item Depressive symptoms in response to a stressor (health psychology, N=182 edges)
        \item Social norms and obesity prevalence (health behavior, N=90 edges)
        \item Older persons emergency department visits (healthcare systems, N=1119 edges)
    \end{enumerate}
    \item \textbf{Two independent runs per CLD} = 6 primary analyses
    \item \textbf{Total measurements:} 32 (including multiple runs for robustness)
    \item \textbf{Approach:} Meta-analysis with one-sample t-tests for statistical validation
\end{itemize}

\section{Performance Across All Metrics}

Table~\ref{tab:rq2_aggregate_all_metrics} presents the performance of all nine context-insensitive metrics across the three CLDs. Metrics are sorted by mean Area Under the Curve (AUC), with statistical significance tested against chance (AUC=0.5) using one-sample t-tests.

\begin{table}[h]
\centering
\small
\caption{CI Metrics Performance: Multi-CLD Meta-Analysis (N=3 CLDs, 6 runs)}
\label{tab:rq2_aggregate_all_metrics}
\begin{tabular}{lcccccc}
\toprule
\textbf{Metric} & \textbf{Mean AUC} & \textbf{95\% CI} & \textbf{Mean r} & \textbf{p-value} & \textbf{Std Dev} & \textbf{Sig.} \\
\midrule
Perplexity                & 0.679 & [0.650, 0.709] & +0.200 & <0.001 & 0.076 & *** \\
Max Window Entropy        & 0.667 & [0.639, 0.695] & +0.211 & <0.001 & 0.072 & *** \\
Min Probability           & 0.641 & [0.617, 0.664] & -0.162 & <0.001 & 0.060 & *** \\
Judge Perplexity          & 0.601 & [0.524, 0.677] & -0.034 & 0.020  & 0.161 & *   \\
Judge Min Probability     & 0.575 & [0.525, 0.624] & -0.031 & 0.009  & 0.104 & **  \\
Judge Max Window Entropy  & 0.559 & [0.509, 0.609] & +0.091 & 0.035  & 0.105 & *   \\
Aggregate Score           & 0.431 & [0.370, 0.493] & +0.052 & 0.043  & 0.129 & *   \\
Judge Cosine Similarity   & 0.265 & [0.232, 0.298] & +0.260 & <0.001 & 0.069 & *** \\
Cosine Similarity         & 0.230 & [0.214, 0.245] & +0.256 & <0.001 & 0.040 & *** \\
\bottomrule
\end{tabular}
\end{table}

\textit{Note: Significance levels indicate whether metrics perform above chance (AUC > 0.5): *** p<0.001, ** p<0.01, * p<0.05, ns = not significant}

\section{Statistical Significance Summary}

\subsection{Correlation with Hallucinations}

\textbf{7/9 metrics} show statistically significant correlation with hallucination status (p<0.05):
\begin{itemize}
    \item \textbf{Significant:} Perplexity, Max Window Entropy, Min Probability, Cosine Similarity, Judge Perplexity, Judge Max Window Entropy, Judge Cosine Similarity
    \item \textbf{Not significant:} Judge Min Probability (p=0.198), Aggregate Score (p=0.058)
\end{itemize}

\subsection{Classification Performance}

\textbf{All 9 metrics} perform significantly above chance (AUC > 0.5, all p<0.05), demonstrating that context-insensitive signals contain measurable information about hallucinations.

\subsection{Cross-CLD Stability}

\textbf{5/9 metrics} exhibit low variance across CLDs (std < 0.1), indicating consistent performance:
\begin{itemize}
    \item \textbf{High stability (std < 0.1):} Cosine Similarity, Judge Cosine Similarity, Max Window Entropy, Min Probability, Perplexity
    \item \textbf{High variance (std > 0.15):} Judge Perplexity (std=0.161)
\end{itemize}

\subsection{Practical Performance}

Maximum F1 score achieved: \textbf{0.364} (perplexity at threshold=1.334)

This performance is:
\begin{itemize}
    \item \textbf{Insufficient} for automated filtering without human oversight
    \item \textbf{Viable} for triage and flagging high-risk edges for manual review
\end{itemize}

\section{The Cosine Similarity Paradox}

A critical methodological finding emerged: both cosine similarity metrics show \textit{positive} correlation with hallucinations, contradicting theoretical expectations.

\subsection{Expected vs. Observed Behavior}

\textbf{Expected:} Higher cosine similarity between generated motivation and citation text should indicate fewer hallucinations (negative correlation), as the LLM's narrative aligns with source material.

\textbf{Observed:} Higher similarity correlates with \textit{more} hallucinations (r=+0.256, p<0.001), suggesting a paradoxical relationship.

Table~\ref{tab:cosine_paradox} presents the paradoxical correlations for all similarity-based metrics.

\begin{table}[h]
\centering
\caption{Cosine Similarity Paradox: Expected vs. Observed Correlations}
\label{tab:cosine_paradox}
\begin{tabular}{lcccl}
\toprule
\textbf{Metric} & \textbf{Expected Sign} & \textbf{Observed r} & \textbf{p-value} & \textbf{Interpretation} \\
\midrule
Cosine Similarity (Gen)   & Negative & +0.256 & <0.001 & Paradoxical \\
Cosine Similarity (Judge) & Negative & +0.260 & <0.001 & Paradoxical \\
Aggregate Judge Score     & Negative & +0.052 & 0.058  & Weak/NS \\
\bottomrule
\end{tabular}
\end{table}

\subsection{Hypotheses for the Paradox}

Two hypotheses may explain this counterintuitive finding:

\subsubsection{Hypothesis 1: Ground Truth Incompleteness}

The validation CLDs may be \textit{incomplete} rather than representing absolute ground truth. Edges labeled as false positives (FP) might represent valid causal relationships that are simply absent from the validation CLD. In this scenario:

\begin{itemize}
    \item High cosine similarity indicates genuine causal relationships well-supported by literature
    \item These edges are incorrectly labeled as "hallucinations" due to validation CLD gaps
    \item The positive correlation reflects ground truth quality issues, not metric failure
\end{itemize}

\textbf{Supporting evidence:} In prior analysis (not shown), correlations exhibited sign-flip when using judge aggregate score as ground truth instead of FP labels, suggesting the FP labels themselves may be unreliable.

\subsubsection{Hypothesis 2: Max-Aggregation Flaw}

Cosine similarity computation uses max() aggregation over citation chunks:

\begin{verbatim}
gen_cosine_sim = max(sims) if sims else None
\end{verbatim}

This approach rewards finding \textit{one} topically-similar chunk while ignoring whether the citation supports the \textit{causal mechanism}. The system can achieve high similarity scores by:

\begin{itemize}
    \item Cherry-picking chunks that discuss relevant topics
    \item Ignoring whether these chunks establish causation
    \item Combining multiple partial-match citations into a seemingly well-supported claim
\end{itemize}

\textbf{Important note:} Citations are retrieved \textit{post-hoc} via Brave search based on the generated motivation text. The LLM cannot deliberately generate high-similarity narratives to match citations, as retrieval occurs after generation completes. This rules out deliberate confabulation strategies.

\subsection{Implications}

The paradox suggests:
\begin{enumerate}
    \item Cosine similarity, as currently implemented, may be \textit{inversely} useful: lower similarity might indicate more reliable edges
    \item Ground truth validation requires expert review, not just CLD completeness checks
    \item Alternative similarity aggregation methods (mean, weighted by chunk position) warrant investigation
\end{enumerate}

\section{Answer to Research Question 2}

\subsection{Research Question}

Can context-insensitive metrics predict hallucinations in LLM-generated causal discovery outputs?

\subsection{Answer: PARTIALLY - with Important Caveats}

\subsubsection{What Works}

\begin{itemize}
    \item \textbf{Uncertainty-based metrics} (perplexity, max window entropy) show moderate, statistically significant predictive power:
    \begin{itemize}
        \item Perplexity: AUC=0.679 (95\% CI [0.650, 0.709])
        \item Max Window Entropy: AUC=0.667 (95\% CI [0.639, 0.695])
    \end{itemize}
    \item \textbf{Statistical robustness:} All metrics show significant correlation or classification performance (p<0.05)
    \item \textbf{Cross-CLD generalization:} Top metrics exhibit low variance (std < 0.1), indicating consistent performance across different health contexts
    \item \textbf{Measurable signal:} Even the weakest predictor (cosine similarity, AUC=0.230) performs significantly above chance
\end{itemize}

\subsubsection{What Doesn't Work}

\begin{itemize}
    \item \textbf{Similarity-based metrics} exhibit paradoxical behavior, likely due to ground truth quality issues or max-aggregation flaws
    \item \textbf{Practical performance} insufficient for automated filtering:
    \begin{itemize}
        \item Maximum F1 score: 0.364 (perplexity)
        \item Precision-recall tradeoff remains unfavorable
    \end{itemize}
    \item \textbf{Judge-based metrics} generally underperform generator-based metrics:
    \begin{itemize}
        \item Aggregate Score: AUC=0.431 (barely above chance)
        \item Judge metrics show higher variance across CLDs
    \end{itemize}
\end{itemize}

\subsubsection{Practical Implications}

\textbf{Recommended use cases:}
\begin{itemize}
    \item \textbf{Triage/flagging:} Use perplexity or max window entropy to prioritize edges for human review
    \item \textbf{Multi-metric ensemble:} Combine multiple CI metrics rather than relying on single thresholds
    \item \textbf{Risk stratification:} Create low/medium/high risk categories rather than binary accept/reject
\end{itemize}

\textbf{Not recommended:}
\begin{itemize}
    \item Automated filtering without human oversight
    \item Sole reliance on similarity-based metrics
    \item Single-metric decision boundaries
\end{itemize}

\section{Limitations and Future Work}

\subsection{Limitations}

\begin{enumerate}
    \item \textbf{Domain limitation:} Validation limited to 3 CLDs within health domain. Generalization to other domains (economics, ecology, social science) remains unknown.
    
    \item \textbf{Ground truth quality:} Validation CLDs may be incomplete. The cosine similarity paradox suggests systematic issues with FP labeling that affect all metrics.
    
    \item \textbf{Max-aggregation method:} Current cosine similarity implementation may be fundamentally flawed for causal assessment. Alternative aggregation strategies unexplored.
    
    \item \textbf{Sample size:} Only 3 CLDs limits statistical power for subgroup analyses and domain-specific effects.
    
    \item \textbf{Judge underperformance:} Judge-based metrics consistently underperform generator metrics, suggesting the judge LLM may evaluate citation quality rather than causal correctness.
\end{enumerate}

\subsection{Future Work}

\begin{enumerate}
    \item \textbf{Domain expansion:} Validate metrics across non-health domains to assess true generalizability
    
    \item \textbf{Ground truth validation:} Expert review of high-cosine-similarity FP edges to validate or refute incompleteness hypothesis
    
    \item \textbf{Alternative similarity methods:}
    \begin{itemize}
        \item Mean aggregation over chunks instead of max
        \item Weighted aggregation by chunk relevance
        \item Causal-specific similarity metrics that weight mechanism descriptions
    \end{itemize}
    
    \item \textbf{Context-sensitive metrics:} Develop metrics that explicitly evaluate causal mechanism support, not just topic overlap
    
    \item \textbf{Ensemble methods:} Investigate optimal combinations of CI metrics for improved F1 scores
    
    \item \textbf{Judge improvement:} Retrain or prompt-engineer judge LLM to focus on causal correctness rather than citation formatting
\end{enumerate}

\section{Conclusion}

Context-insensitive metrics provide measurable, statistically significant signals for hallucination detection in LLM-generated causal discovery, with uncertainty-based metrics (perplexity, entropy) demonstrating moderate predictive power (AUC $\sim$0.67) that generalizes across multiple CLDs. However, practical performance remains insufficient for automated filtering (F1<0.5), and similarity-based metrics exhibit paradoxical behavior suggesting fundamental issues with either ground truth quality or metric implementation. The most viable current application is human-in-the-loop triage, where CI metrics flag high-risk edges for manual review rather than serving as autonomous gatekeepers.



